% Gauss Werke Band 11, Abt. 1 - Pages 251-300 (original pp. 249-298)
% Teil 6: Theoretische und Praktische Astronomie
\documentclass[12pt,a4paper]{article}
\usepackage[utf8]{inputenc}
\usepackage[T1]{fontenc}
\usepackage[ngerman,latin]{babel}
\usepackage{amsmath,amssymb,amsthm}
\usepackage{geometry}
\usepackage{hyperref}
\usepackage{enumitem}
\usepackage{mathrsfs}
\usepackage{longtable}
\usepackage{booktabs}
\usepackage{microtype}
\usepackage{array}
\usepackage{multirow}
\usepackage{tabularx}
\usepackage{mathtools}
\usepackage{oldgerm}
\usepackage{textcomp}
\newcommand{\degree}{\ensuremath{^\circ}}
\geometry{margin=2.5cm}

\theoremstyle{plain}
\newtheorem{satz}{Satz}[section]
\newtheorem{theorem}{Theorem}[section]
\newtheorem{lemma}{Lemma}[section]
\newtheorem{hilfssatz}{Hilfssatz}[section]
\newtheorem{korollar}{Korollar}[section]

\theoremstyle{definition}
\newtheorem{definition}{Definition}[section]

\theoremstyle{remark}
\newtheorem{bemerkung}{Bemerkung}[section]
\newtheorem{beispiel}{Beispiel}[section]

\title{Teil 6}
\author{Carl Friedrich Gau\ss}
\date{}

\begin{document}
\maketitle

% ============================================
% SECTION: ZUR BAHNBESTIMMUNG DER CERES (Schluss)
% ============================================
\section*{Zur Bahnbestimmung der Ceres}

so wird diese Gleichung\footnote{Vgl.\ Theoria motus, art.~120--122.}
\begin{align*}
6639 &= 7785x' + 4598y,\\
9945 &= 11750x' + 8407{,}5y,
\end{align*}
[wo $v = \frac{\Delta Q}{677}$ und $\frac{\Delta i}{\pi r}$ gesetzt ist.

Die Auflösung ergibt]
\begin{align*}
\log x' &= 9{,}9460893\quad (\log y = 8{,}7125n),\quad\text{(also}\\
\Delta Q &= 385{,}99\quad [\Delta i = -9{,}08].
\end{align*}
[Es ergibt sich]
\begin{center}
\begin{tabular}{c|c}
$80\degree55'18''$ & $10\degree36'30''$ \\
$\phantom{0}6\phantom{\degree}\,\,26$ & $-9{,}08$ \\
\hline
$\Omega = \,]\,81\degree\,14'\,4$ & $[i = \,]\,10\degree\,36'\,20{,}92$.
\end{tabular}
\end{center}

(S.~12 u.~13.) [Mit diesen Werten als] Vierte Hypothese [wurde die] Berechnung der Bahn [von neuem ausgeführt und die III.~Elemente gefunden\footnote{Vgl.\ Werke~VI, S.~201.}.]

\subsection*{Bemerkungen}

Die ältesten im Nachlaß erhaltenen Gaussschen Rechnungen und Notizen zur Bahnbestimmung, insbesondere für Ceres, finden sich in den Schedae Ag und Ah, sowie im Handbuch Bb, und zwar in den ersten beiden vorzugsweise die numerischen Rechnungen, im letzteren dagegen nur Formelentwicklungen; sie stammen aus dem November 1801, sodaß also die allerersten Untersuchungen aus dem September und Oktober (vergl.\ das Tagebuch No.~119 u.~121, Werke~XI, S.~561, 563) verloren gegangen sind. Man wird aber in der Notiz [1], oben S.~221 ff., die auf den ersten Seiten des Handbuchs Bb unter der Überschrift »Erste Abhandlung« steht, einen Niederschlag der in der Tagebuchnotiz No.~121 erwähnten »permultae formulae« erblicken dürfen.

Die Notiz [II], oben S.~232 ff., steht in der Scheda Ag. Die Zahlenbeispiele bilden die fortlaufende Rechnung zur Ableitung der I.~Elemente der Ceres, die Gauss an v.~Zach schickte (Werke~VI, S.~200~f.). Die vorausgesetzten Werte von $i$ und $\Omega$ sind vielleicht einer noch früheren, verloren gegangenen Originalrechnung\footnote{Gauss sagt in der Einleitung zur Theoria motus (Werke~VII, 1906, S.~8), daß er die erste Bahnbestimmung der Ceres im Oktober 1801 gemacht habe.} entnommen; bemerkenswert ist, daß Gauss von der hier angewandten »Dritten Verbesserungsmethode« (siehe S.~229) später in der Theoria motus, art.~2 (vergl.\ auch art.~126) sagt, daß sie »orbitae dimensiones eruendi magnam praecisionem nunquam admittet, nisi tria loca heliocentrica intervallis considerabilibus ab invicem distent«.

\vspace{0.5em}
\noindent\textsc{Die Notiz [III]}, oben S.~241 ff., steht in der Scheda Ah; es ist dies das Original der Berechnung der Ceres-Elemente III. durch Verbesserung der Elemente II. Die numerischen Rechnungen sind beim Abdruck gekürzt worden, insbesondere wurden die bei drei Beobachtungen gleichlautenden Rechnungen nur bei der ersten Beobachtung ausführlich abgedruckt.

Wie Gauss nicht nur in der Einleitung zur Theoria motus und in einem Briefe an \textsc{Olbers} vom 6.~August 1802 sagt, sondern wie auch aus der Anmerkung hervorgeht, mit der von \textsc{Lindenau} die Herausgabe des ihm von Gauss übergebenen Manuskripts \textit{Summarische Übersicht der zur Bestimmung der Bahnen der beiden neuen Hauptplaneten angewandten Methoden} (Monatl.\ Corresp., Sept.\ 1809 und Werke, Bd.~VI, S.~148) begleitet hat, unterscheiden sich die Methoden der Theoria motus sehr wesentlich von den ältesten Gaussschen Rechnungsweisen; die ersteren sind nach den Tagebuchnotizen No.~125--129 (Werke~XI, S.~564, 565) auch erst in den Jahren 1805/06 entstanden. Doch erwähnt Gauss auch einige der letzteren gelegentlich in der Theoria motus; auf die betreffenden Parallelstellen ist oben durch Fußnoten hingewiesen. Die erwähnte \textit{Summarische Übersicht} hatte Gauss schon im Jahre 1802 zusammengestellt und am 6.~August an Olbers gesandt (siehe \textsc{Wilhelm Olbers}, Sein Leben und seine Werke~II, 1, 1900, S.~65).

In den Artikeln 124--129 der Theoria motus bespricht Gauss zehn verschiedene Methoden der Bahnbestimmung; man wird bemerken, daß die unter (I.) abgedruckten Notizen [3]--[6] sich an die erste dieser Methoden, die Notiz [7] an die siebente (Theoria motus, art.~129, No.~II) und die Notiz [8] an die achte (Theoria motus, art.~129, No.~III) anschließt; doch fehlt bei der Notiz [8] der Hinweis, daß die Interpolation auf die verbesserten Werte durch Vergleichung der berechneten Zwischenzeiten mit den beobachteten erfolgen soll. Die Notiz [II.] schließt sich ebenfalls an die achte Methode und die Notiz [III] an die zweite (Theoria motus, art.~126) an, obgleich die numerische Durchführung der letzteren in den Notizen abweicht.

\subsection*{Bemerkungen (Fortsetzung)}

Führt man die Differentiation auf den linken Seiten der Gleichungen (3)--(5) aus, multipliziert sie zur Elimination von $\delta e$ und $\delta\varepsilon$ der Reihe nach mit
\[
0,\quad \frac{dF}{da}\sin\alpha - \frac{dF}{d\delta}\cos\alpha,\quad
\frac{dF}{da}\cos\alpha + \frac{dF}{d\delta}\sin\alpha,\quad \frac{dF}{d\tau}
\]
und addiert, so wird
\begin{equation}
\delta\alpha = \frac{AB}{a'}\Bigl(\frac{\delta\Theta}{\Theta} + \frac{\delta p}{p}\Bigr)
\tan\psi\cos(V-\alpha) + \frac{\delta R}{R}\sin(V-\alpha)
\tag{a}
\end{equation}
und dies ist die Gleichung (a) der Nachlaßnotiz, wo nur $A=1$ gesetzt ist. Man vergleiche hierzu den Brief von Olbers an Gauss vom 11.~September 1802 und \textsc{Laplace}, \textit{Mécanique céleste}, Première Partie, Livre~II, No.~31 (Band~I, S.~207 der Originalausgabe von 1799).

Um die Gleichung (b) abzuleiten, multipliziert man in ähnlicher Weise die Gleichungen (3)--(5) zur Elimination von $\delta e$ und $\delta\varepsilon$ der Reihe nach mit $-\delta\Theta\sin V$, $\Theta\cos V$, $\sin(V-\alpha)$ und addiert.

Die Gleichungen (a) und (b) lassen sich auch aus I und II durch den Grenzübergang $\lim\Delta U = 0$ ableiten, wobei noch zu beweisen ist, daß\footnote{Der Beweis ist in der That nicht ganz einfach.} die in I und II vernachlässigten Glieder beim Grenzübergang verschwinden.

Die Gleichung (c) erhält man durch Multiplikation der Gleichung (3) bezw.\ (4) mit $-\sin\alpha$ bezw.\ $\cos\alpha$ und Addition und endlich die Gleichung (d) durch Multiplikation der Gleichungen (3)--(5) mit $-\Theta\cos\alpha$, $-\Theta\sin\alpha$, $+1$ und Addition.

Die Gleichungen für $\log A'$ und $\log A''$, S.~223, beruhen auf der Voraussetzung, daß die Logarithmen der Abstände von der Erde sich proportional den Zwischenzeiten ändern, daß also
\[
\frac{\log A'' - \log A'}{\tau'' - \tau'} = \frac{\log A' - \log A}{\tau' - \tau}.
\]

Mit der auf S.~229 bei Gelegenheit der »zweiten Verbesserungsmethode« erwähnten »Aufgabe« hat sich Gauss öfters beschäftigt; es finden sich darüber Notizen im Handbuch Bb, §.~15 und in der Scheda Ah, S.~5 (mit Beispielen) und Gauss hat diesen Teil seiner Untersuchungen in der Monatl.\ Corresp.\ Juni 1802 (Bd.~VI, S.~87) veröffentlicht; vergl.\ auch Theoria motus, art.~74.

Vielleicht ist es erwünscht, die Ableitung der Formel für $p$ in der Fußnote S.~225 aus der vorhergehenden Gleichung hier zu geben, da sie nicht sofort zu übersehen ist. Dabei ergibt sich der genaue Wert:
\[
\frac{599}{5120} = \frac{11}{26}\text{ [fehlt, eben]}
\]
während der genaue Wert ist:
\[
\frac{599}{5120}.
\]

\vspace{0.3em}
\noindent Außer den oben beim Text angegebenen wurden noch einige weitere kleinere Schreibfehler beim Abdruck verbessert.

Es schien angebracht, im vorstehenden die ältesten Urkunden der Arbeiten von Gauss auf diesem Gebiete zum Abdruck zu bringen, wenn auch viele Einzelheiten aus der \textit{Summarischen Übersicht} und der Theoria motus entnommen werden können.

\begin{flushright}
\textsc{Brendel.}
\end{flushright}

\newpage
% ============================================
% SECTION: DIE GAUSSSCHE METHODE ZUR BERECHNUNG EINER KREISBAHN
% ============================================
\section*{Die Gaußsche Methode zur Berechnung einer Kreisbahn}
\footnote{Der Verfasser [Klinkerfues] verdankt diese Methode den mündlichen Mitteilungen von Gauss, bei Gelegenheit eines Auftrages zur Untersuchung der Bahn der Eunomia.}

[Aus \textsc{W.~Klinkerfues}, \textit{Theoretische Astronomie}, 1.~Aufl., Braunschweig 1871, S.~63; 3.~Aufl., herausgegeben von H.~Buchholz, Braunschweig 1912, S.~237.]

Das Verfahren von Gauss, an zwei vollständige Beobachtungen eines Planeten eine Kreisbahn anzuschließen, gewährt für die Ekliptik eine sehr bequeme Rechnung.

Wir beginnen damit, uns die Lage der Bahnebene gegen die Ekliptik, sowie die Bedingungen der Aufgabe \ldots\ in [der beigefügten] Figur zu versinnlichen.

\subsection*{}

Es sei $EC$ der die Ekliptik vorstellende größte Kreis der Himmelskugel, $T$ und $T'$ seien die Orter der Erde auf derselben, d.\,h.\ diejenigen Punkte der Sphäre, auf welche die Radienvectoren der Erde in den beiden Beobachtungen gerichtet sind. Durch jede vollständige Beobachtung wird die Richtung oder eine Gerade, auf welcher der Planet sich zur Zeit der Beobachtung befinden muß, völlig bekannt, demnach auch die Lage der Ebene, welche durch diese Gerade und den Radiusvector der Erde in derselben Beobachtung gelegt werden kann. Diese Ebenen für die erste und die zweite Beobachtung werden in der Figur durch die größten Kreise $TD$ und $T'D$, welche sich in $D$ schneiden, vorgestellt. $QPP'O$ sei die Bahnebene des Planeten, worin $Q$ der aufsteigende Knoten in der Ekliptik, $P$ und $P'$ die heliocentrischen Orter des Planeten in beiden Beobachtungen. Die geocentrischen Orter, in der Figur mit $\Pi$ und $\Pi'$ bezeichnet, fallen offenbar ebenfalls in die größten Kreise $TD$ und $T'D$. Man kann ferner noch leicht bemerken, daß der Bogen $TP$ nichts Anderes ist, als der Winkel, welchen der Radiusvector des Planeten mit dem Radiusvector der Erde in der ersten Beobachtung bildet, $T'P'$ dasselbe für die zweite Beobachtung. Diese beiden Bogen sind nicht von vornherein gegeben, dagegen die Bogen $T\Pi$ und $T'\Pi'$, d.\,h.\ die Winkel, welche an der Erde die Beobachtungsrichtungen mit den Verlängerungen der Radienvectoren der Erde in beiden Beobachtungen bilden; diese

\begin{align*}
T\Pi &= \lambda - L, \\
T'\Pi' &= \lambda' - L', \\
\Pi A &= \beta, \\
\Pi' A' &= \beta',
\end{align*}
also nach den Formeln rechtwinkliger sphärischer Dreiecke
\begin{align}
\sin\beta \tan\gamma &= \sin(\lambda - L), \tag{1}\\
\sin\beta' \tan\gamma' &= \sin(\lambda' - L'). \tag{2}
\end{align}

Die Formel für $\chi$:
\[
\cos\chi = \cos\beta\cos(\lambda - L)
\]
kann zur Controle der Rechnung dienen. Für einen eben entdeckten noch in der Nähe seiner Opposition befindlichen Planeten, pflegt der Winkel $\gamma$ nicht groß zu sein, weshalb seine Bestimmung durch (1) und (2) der anderen vorzuziehen ist.

Bei der zweiten Beobachtung hat man analog:
\begin{align}
\tan\gamma' &= \frac{\sin(\lambda' - L')}{\sin\beta'}, \tag{3}\\
\sin\gamma' &= \frac{\cos\beta'}{\text{[etwas]}}. \tag{4}
\end{align}

Nach der Berechnung von $\gamma$ und $\gamma'$ kann man nun sofort zur Auflösung des Dreiecks $TDT'$, d.\,h.\ zur Bestimmung der beiden Seiten $TD$, $T'D$ und des Winkels $TDT'$ oder $\varepsilon$, aus den Winkeln $\gamma$, $(180\degree - \gamma')$ und der dazwischenliegenden Seite $TT'$ oder $L' - L$ schreiten. Die Gaußschen Dreiecksformeln liefern hier die Gleichungen:
\begin{align}
\sin\tfrac{1}{2}\varepsilon \cdot \sin\tfrac{1}{2}(TD + T'D) &= \sin\tfrac{1}{2}(L' - L) \sin\tfrac{1}{2}(\gamma + \gamma'), \tag{5}\\
\sin\tfrac{1}{2}\varepsilon \cdot \cos\tfrac{1}{2}(TD + T'D) &= \cos\tfrac{1}{2}(L' - L) \sin\tfrac{1}{2}(\gamma' - \gamma), \tag{6}\\
\cos\tfrac{1}{2}\varepsilon \cdot \sin\tfrac{1}{2}(TD - T'D) &= \sin\tfrac{1}{2}(L' - L) \cos\tfrac{1}{2}(\gamma' + \gamma), \tag{7}\\
\cos\tfrac{1}{2}\varepsilon \cdot \cos\tfrac{1}{2}(TD - T'D) &= \cos\tfrac{1}{2}(L' - L) \cos\tfrac{1}{2}(\gamma' - \gamma). \tag{8}
\end{align}

Hat man diese Vorbereitungsrechnungen beendigt, so erhält man für jede Hypothese über den Werth von $r$ oder den Abstand des Planeten von der Sonne sofort die Bogen $DP$ und $DP'$, aus denen man dann in Verbindung mit dem von ihnen eingeschlossenen Winkel $\varepsilon$ die Seite $PP'$ des Dreiecks $PDP'$, d.\,h.\ den heliocentrischen Bogen $PP'$, welcher, wenn die Hypothese richtig ist, gleich $\alpha\sqrt[3]{\frac{\Delta t}{2}}$ sein soll, bestimmen kann. Denn ein Blick auf die Figur zeigt, daß:
\begin{align}
DP &= DT - PT = DT - (\gamma - \chi), \tag{9}\\
DP' &= DT' - P'T' = DT' - (\gamma' - \chi'), \tag{10}
\end{align}
während
\[
\sin\chi = \frac{R}{r}\sin\gamma,\quad \sin\chi' = \frac{R'}{r}\sin\gamma',
\]
wobei $R$ und $R'$ die Radienvectoren der Erde vorstellen.

Setzen wir der Kürze halber:
\[
\text{Winkel }DPO = \eta,\quad \text{Winkel }DP'O = \eta',
\]
so haben wir wieder nach den Gaussschen trigonometrischen Formeln:
\begin{align}
\sin\tfrac{1}{2}PP' \sin\tfrac{1}{2}(\eta' + \eta) &= \sin\tfrac{1}{2}\varepsilon \sin\tfrac{1}{2}(DP + DP'), \tag{11}\\
\cos\tfrac{1}{2}PP' \sin\tfrac{1}{2}(\eta' - \eta) &= \sin\tfrac{1}{2}\varepsilon \cos\tfrac{1}{2}(DP + DP'), \tag{12}\\
\sin\tfrac{1}{2}PP' \cos\tfrac{1}{2}(\eta' + \eta) &= \cos\tfrac{1}{2}\varepsilon \sin\tfrac{1}{2}(DP - DP'), \tag{13}\\
\cos\tfrac{1}{2}PP' \cos\tfrac{1}{2}(\eta' - \eta) &= \cos\tfrac{1}{2}\varepsilon \cos\tfrac{1}{2}(DP - DP'), \tag{14}
\end{align}
aus denen $PP'$, $\eta'$ und $\eta$ sich ergeben. Die beiden letzteren Größen $\eta$ und $\eta'$ gewinnen erst Interesse, wenn man den definitiven Werth von $r$ gefunden hat, weshalb es vortheilhaft ist, dieselben aus den Versuchen zu eliminiren. Man addire zu dem Zwecke die Gleichungen (11) und (13), nachdem man auf beiden Seiten in das Quadrat erhoben hat, wodurch sich ergibt:
\begin{equation}
\sin^2\tfrac{1}{2}PP' = \sin^2\tfrac{1}{2}\varepsilon \sin^2\tfrac{1}{2}(DP + DP') + \cos^2\tfrac{1}{2}\varepsilon \sin^2\tfrac{1}{2}(DP - DP'), \tag{15}
\end{equation}
(während kaum nöthig scheint, zu bemerken, daß sich der Exponent 2 nicht auf die in Klammern stehenden Bogen, sondern auf das Quadriren des Sinus bezieht).

Hat man die definitive Lösung gefunden, so rechnet man nach den Formeln (9) und (12), so daß auch $\eta$ und $\eta'$ bekannt werden. Wieder ein Blick auf die Figur zeigt, daß die Bestimmung der Elemente $i$ und $\Omega$ aus den Winkeln $\eta$, $\gamma$, $\chi$, sowie aus dem Längenunterschied $L - \Omega$ durch Auflösung des Dreiecks $PQ\Pi$ finden sind. In diesem ist offenbar:
\begin{align*}
\text{Winkel }P\Pi T &= i = \text{Neigung der Bahnebene zur Ekliptik},\\
QP &= u,\\
Q\Pi &= L - \Omega,
\end{align*}
und da auch Winkel $QP\Pi = \chi$, $P\Pi Q = 180\degree - \gamma$, $PT\Pi = \gamma - \chi$, so ist nach den Dreiecksformeln:
\begin{align}
\sin\tfrac{1}{2}i \sin\tfrac{1}{2}[u + (L - \Omega)] &= \sin\tfrac{1}{2}(\gamma - \chi) \sin\tfrac{1}{2}(\gamma + \eta), \tag{16}\\
\sin\tfrac{1}{2}i \cos\tfrac{1}{2}[u + (L - \Omega)] &= \cos\tfrac{1}{2}(\gamma - \chi) \sin\tfrac{1}{2}(\gamma + \eta), \tag{17}\\
\cos\tfrac{1}{2}i \sin\tfrac{1}{2}[u - (L - \Omega)] &= \sin\tfrac{1}{2}(\gamma - \chi) \cos\tfrac{1}{2}(\gamma + \eta), \tag{18}\\
\cos\tfrac{1}{2}i \cos\tfrac{1}{2}[u - (L - \Omega)] &= \cos\tfrac{1}{2}(\gamma - \chi) \cos\tfrac{1}{2}(\gamma - \eta). \tag{19}
\end{align}

Zur Controle der Rechnung dient die Übereinstimmung des $i$ oder der Neigung, wie sie aus (16) und (17) und aus $\sin\tfrac{1}{2}i$ folgt, mit der aus (18) und (19) und aus $\cos\tfrac{1}{2}i$.

Zu weiterer Sicherung der Richtigkeit kann man noch das Dreieck $QP'T'$ zur Bestimmung der Elemente verwenden.

Man hat darnach auch:
\begin{align}
\sin\tfrac{1}{2}i \sin\tfrac{1}{2}[u' + (L' - \Omega)] &= \sin\tfrac{1}{2}(\gamma' - \chi') \sin\tfrac{1}{2}(\gamma' + \eta'), \tag{20}\\
\sin\tfrac{1}{2}i \cos\tfrac{1}{2}[u' + (L' - \Omega)] &= \cos\tfrac{1}{2}(\gamma' - \chi') \sin\tfrac{1}{2}(\gamma' - \eta'), \tag{21}\\
\cos\tfrac{1}{2}i \sin\tfrac{1}{2}[u' - (L' - \Omega)] &= \sin\tfrac{1}{2}(\gamma' - \chi') \cos\tfrac{1}{2}(\gamma' + \eta'), \tag{22}\\
\cos\tfrac{1}{2}i \cos\tfrac{1}{2}[u' - (L' - \Omega)] &= \cos\tfrac{1}{2}(\gamma' - \chi') \cos\tfrac{1}{2}(\gamma' - \eta'). \tag{23}
\end{align}
Das Argument der Breite in der zweiten Beobachtung, $u'$, muß noch der Bedingung genügen, daß $u' - u = PP' = \alpha\sqrt[3]{\frac{\Delta t}{2}}$ wird, wie es immer der Fall sein wird, wenn die Rechnung vorher in allen Theilen scharf geführt worden ist.

\newpage
% ============================================
% SECTION: PRAKTISCHE UND SPHARISCHE ASTRONOMIE
% ============================================
\section*{Praktische und Sphärische Astronomie}
\subsection*{Nachträge zu Band VI}

\newpage
% ============================================
% SUBSECTION: ASTRONOMISCHE BEOBACHTUNGEN AUS DER BRAUNSCHWEIGER ZEIT
% ============================================
\subsection*{[Astronomische Beobachtungen aus der Braunschweiger Zeit]}

\subsubsection*{Veröffentlichungen und Briefwechsel}

\paragraph{(I) [Längenbestimmung von Braunschweig, Helmstedt und Wolfenbüttel].}
\textit{[Monatliche Correspondenz, Band X, S.~302--307, 1804 October. Aus v.~Zach's Bericht über die Kgl.\ Preussische trigonometrische Aufnahme von Thüringen u.\,s.\,w.]}

\subparagraph{1.~Braunschweig.}
\textit{[Der geheime Rath Freyherr von Ende und Dr.~Gauss in Braunschweig hatten die Gefälligkeit, unsere Brockensignale in Braunschweig, Helmstedt und Wolfenbüttel zu beobachten. Der geh.\ R.\ v.\ E.\ hatte seinen eigenen Arnold'schen Chronometer, und dem Dr.\ Gauss überschickte ich einen dergleichen zu diesem Behufe. Ersterer beobachtete in seiner Wohnung in der Steinstraße, letzterer an der Südseite der Stadt auf dem Garten des Kaufmanns Köpp, etwa 400 Par.\ Fuß westlich, und 3500 Fuß südlich vom Andreas-Thurm. Des Freyh.\ v.\ Ende Wohnung ist 135' östlich von Köpp's Garten. Da die Witterung nicht die günstigste war, so mußte man bisweilen zu einzelnen Höhen seine Zuflucht nehmen, um die Zeitbestimmung zu erhalten, worin auch wahrscheinlich die kleinen Differenzen zu suchen sind, welche sich in der Bestimmung der Länge von Braunschweig zeigen.\footnote{Zach's eigene Beobachtungen finden sich nicht hier, sondern Monatl.\ Corresp., Bd.~X, S.~304.}]]}

\begin{center}
\small
\begin{tabular}{|r|c|c|c|}
\hline
\multicolumn{4}{|c|}{\textit{1803}} \\
\hline
& \textit{Mittlere Zeit} & \textit{Mittlere Zeit} & \textit{Länge in Zeit} \\
& \textit{auf Seeberg} & \textit{in Braunschweig} & \textit{westl.\ v.\ Seeberg} \\
\hline
9.~August & $9^h\,11^m\,35^s{,}4$ & $9^h\,10^m\,16^s{,}54$ & $46^s{,}87$ \\
\quad\,21 & \quad 2\,,5 & \quad 20\,,16\,,3 & 46\,,2 \\
\quad\,31 & \quad 2\,,9 & \quad 30\,,16\,,3 & 46\,,7 \\
\quad\,41 & \quad 3\,,3 & \quad 40\,,16\,,2 & 47\,,1 \\
\quad\,51 & \quad 2\,,6 & \quad 50\,,15\,,7 & 46\,,9 \\
\hline
10.\quad 1 & \quad 3\,,4 & 10\,\,0\,,16\,,3 & 47\,,1 \\
\hline
\multicolumn{3}{|l|}{Anzahl d.\ Sign.\ 6} & \textit{Mittel 46{,}78} \\
\hline
\end{tabular}
\end{center}

\begin{center}
\small
\begin{tabular}{|r|c|c|c|}
\hline
15.~August & $8^h\,17^m\,26^s{,}83$ & $9^h\,0^m\,40^s{,}59$ & $45^s{,}84$ \\
\quad\,10 & \quad 27\,,1 & \quad 40\,,9 & 46\,,2 \\
\quad\,20 & \quad 26\,,5 & \quad 19\,,41\,,0 & 45\,,5 \\
\quad\,30 & \quad 26\,,1 & \quad 29\,,41\,,1 & 45\,,0 \\
\quad\,40 & \quad 26\,,4 & \quad 39\,,41\,,5 & 44\,,9 \\
\quad\,50 & \quad 26\,,7 & \quad 49\,,41\,,8 & 44\,,9 \\
\hline
\multicolumn{3}{|l|}{Anzahl d.\ Sign.\ 6} & \textit{Mittel 45{,}32} \\
\hline
\end{tabular}
\end{center}

\begin{center}
\small
\begin{tabular}{|r|c|c|c|}
\hline
17.~August & $8^h\,57^m\,24^s{,}20$ & $8^h\,59^m\,36^s{,}53$ & $47^s{,}57$ \\
\quad\,10 & \quad 24\,,4 & \quad 9\,,36\,,1 & 48\,,3 \\
\quad\,20 & \quad 24\,,5 & \quad 19\,,36\,,3 & 48\,,2 \\
\quad\,30 & \quad 24\,,3 & \quad 29\,,35\,,8 & 48\,,5 \\
\quad\,40 & \quad 24\,,5 & \quad 39\,,36\,,0 & 48\,,5 \\
\quad\,50 & \quad 25\,,0 & \quad 49\,,36\,,4 & 48\,,6 \\
\quad 0 & \quad 25\,,4 & \quad 59\,,36\,,5 & 48\,,9 \\
\hline
\multicolumn{3}{|l|}{Anzahl d.\ Sign.\ 7} & \textit{Mittel 48{,}39} \\
\hline
\end{tabular}
\end{center}

\begin{center}
\small
\begin{tabular}{|r|c|c|c|}
\hline
18.~August & $6^h\,20^m\,18^s{,}50$ & $6^h\,19^m\,28^s{,}83$ & $49^s{,}7$ \\
\quad\,30 & \quad 18\,,4 & \quad 29\,,28\,,3 & 50\,,1 \\
\quad\,40 & \quad 18\,,9 & \quad 39\,,28\,,6 & 50\,,3 \\
\quad\,50 & \quad 18\,,6 & \quad 49\,,28\,,3 & 50\,,3 \\
\quad\,0 & \quad 18\,,6 & \quad 59\,,28\,,3 & 50\,,3 \\
\hline
\multicolumn{3}{|l|}{Anzahl d.\ Sign.\ 5} & \textit{Am 18.~Aug.\ 50{,}14} \\
\hline
\end{tabular}
\end{center}

Am 18.~August 50{,}14\\
\null\quad\quad\quad\,\,\,\,\,\,\,\,17.\quad\quad 48{,}39\\
\null\quad\quad\quad\,\,\,\,\,\,\,\,15.\quad\quad 45{,}32\\
\null\quad\quad\quad\,\,\,\,\,\,\,\,\,9.\quad\quad 46{,}78\\
Anzahl d.\ Sign.\ 24\quad Mittel 47{,}66

\textit{[Handschriftliche Bemerkung von Gauss in seinem Exemplar:]}
\begin{flushleft}
\small
Köppen's Garten\\
\quad 9.\quad 5\quad 50{,}83\\
\quad\quad\quad\quad 47{,}70\\
\quad\quad\quad\quad 51{,}21\\
\quad\quad\quad\quad 46{,}07\\
\quad\quad\quad\quad 48{,}11\\
\quad\quad\quad\quad 51{,}30\\
\quad\quad\quad\quad 49{,}30\\
\quad\quad Mittel 46\quad$49^s{,}14$.
\end{flushleft}

Nach diesen Beobachtungen liegt die Wohnung des Freyh.\ von Ende in Braunschweig $32^m\,47^s{,}34$ von Paris, oder geographische Länge von Ferro $28\degree11'50''$, und Kaufmann Köpp's Garten $32^m\,48^s{,}38$\footnote{Zach gibt in seiner eigenen Ausgabe der Monatl.\ Corresp.\ X, S.~303 durchgangig 5 Secunden kleinere Werte.} oder von Ferro $28\degree12'12{,}5''$.

Auch hier ergibt sich dieselbe Probe wie bey Magdeburg. Dr.~Gauss beobachtete nämlich die Bedeckung desselben Sterns $\varepsilon$ im Widder vom Monde\footnote{$\varepsilon$ Arietis, 1803 Sept.~5.}. Stellen wir die vom Can.\ Daub berechneten wahren Conjunctions-Zeiten abermahls zusammen, und leiten die Längenunterschiede daraus ab, so kommt aus dieser Beobachtung der Längenunterschied von Braunschweig mit Paris aus der Beobachtung

\begin{center}
\begin{tabular}{l@{\quad}l}
\quad\quad\quad von Magdeburg & $32^m\,45^s{,}5$ \\
\quad\quad\quad Wien & 32\quad 50\,,0 \\
\quad\quad\quad Prag & 32\quad 40\,,5 \\
\quad\quad\quad Leipzig & 32\quad 39\,,3 \\
\quad\quad\quad Danzig & 32\quad 50\,,0 \\
\quad\quad\quad \textit{Mittel} & \textit{32}\,\textit{47}\,\textit{55}\,,\textit{4}
\end{tabular}
\end{center}

welche astronomische Bestimmung 2\,\,$^{s}$59\,\,\footnote{Zach in seinen eigenen Ausgaben der Monatl.\ Corresp.\ X, S.~305: 1\,$^s$53.} von der terrestrischen abweicht. Merkwürdig ist, daß die Braunschweiger und Magdeburger Beobachtung dieser Sternbedeckung den Längenunterschied für beyde Städte am genauesten gibt.

Noch eine andere Prüfung dieser Länge machte Dr.~Gauss, indem er auf Köpp's Garten Azimuthe des Brocken mit seinem Sextanten gemessen hatte. Verbunden mit der Polhöhe $52\degree15'30''$ berechnete er daraus die Länge

\begin{center}
\small
\begin{tabular}{|l|c|c|}
\hline
1803 & \quad 6.~Aug. & $52\degree\,15'\,47{,}5''$ \\
& \quad\, 7.\quad\,\,\,\,\,& 55\,,1 \\
& \quad 16.\quad\,\,\,\,& 55\,,2 \\
\hline
& \textit{Mittel} & \textit{52\degree}\,\textit{15'}\,\textit{52{,}6''} \\
\hline
\end{tabular}
\end{center}

Dr.~Gauss hat die Polhöhe auf Köpp's Garten im vorigen Sommer $52\degree15'35''$ gefunden; fast dasselbe folgt aus der Reduction der Polhöhe seiner vorigen Wohnung, die er $52\degree16'5''$ bestimmt hatte und die $30''$ nördlicher liegt. Ich habe im Sept[ember] 1800 diese Breite im Hotel d'Angleterre $52\degree15'43''$ beobachtet (M.~C., II.~Bd., S.~562). Daraus würde die Breite vom Dr.~Gauss etwas kleiner werden, hingegen nach obiger Bestimmung des Freyh.\ v.\ Ende in seiner Wohnung etwas größer.\,\,)

\subparagraph{2.~Helmstedt.}
\textit{[Den 18.~August [1803] verfügten sich der Geh.\ R.\ von Ende und Dr.~Gauss nach Helmstedt, und, nachdem sie daselbst in des Hofraths Pfaff Garten ein Dutzend correspondirende Sonnenhöhen genommen hatten, beobachteten sie hierauf des Abends am 19.~Aug.\ folgende Brocken-Signale:]}

\begin{center}
\small
\begin{tabular}{|r|c|c|c|}
\hline
\multicolumn{4}{|c|}{\textit{1803}} \\
\hline
& \textit{Mittl.\ Zeit} & \textit{Mittl.\ Zeit} & \textit{Länge in Zeit} \\
& \textit{auf Seeberg} & \textit{in Helmstedt} & \textit{östl.\ v.\ Seeberg} \\
\hline
19.~August & $6^h\,0^m\,16^s{,}85$ & $6^h\,1^m\,24^s{,}87$ & $1^m\,8^s{,}02$ \\
\quad\,20 & \quad 16\,,2 & \quad 21\,,24\,,9 & 8\,,7 \\
\quad\,30 & \quad 16\,,2 & \quad 31\,,24\,,6 & 8\,,4 \\
\quad\,40 & \quad 16\,,0 & \quad 41\,,25\,,4 & 9\,,1 \\
\quad\,50 & \quad 16\,,5 & \quad 51\,,25\,,1 & 8\,,6 \\
\quad\, 0 & \quad 16\,,7 & \quad 1\,,25\,,2 & 8\,,5 \\
\hline
\multicolumn{3}{|l|}{9\quad 10\quad\,16\,,6} & 9\quad 11\quad 26 \ldots \\
\hline
\end{tabular}
\end{center}

Den 21.~August wurden correspondirende Höhen im Gasthofe zum Erbprinzen, und hierauf folgende Signale beobachtet:

\begin{center}
\small
\begin{tabular}{|r|c|c|c|}
\hline
\multicolumn{4}{|c|}{\textit{1803}} \\
\hline
& \textit{Mittl.\ Zeit} & \textit{Mittl.\ Zeit} & \textit{Länge in Zeit} \\
& \textit{auf Seeberg} & \textit{in Helmstedt} & \textit{östl.\ v.\ Seeberg} \\
\hline
21.~August & $9^h\,10^m\,5^s{,}86$ & $9^h\,11^m\,14^s{,}85$ & $1^m\,8^s{,}99$ \\
\quad\,20 & \quad 5\,,5 & \quad 21\,,14\,,2 & 8\,,7 \\
\quad\,30 & \quad 5\,,6 & \quad 31\,,14\,,6 & 9\,,0 \\
\quad\,40 & \quad 5\,,8 & \quad 41\,,14\,,3 & 8\,,5 \\
\quad\,50 & \quad 6\,,0 & \quad 51\,,14\,,4 & 8\,,4 \\
\quad\, 0 & \quad 5\,,5 & 10\,\,\,4\,,14\,,5 & 9\,,0 \\
\hline
\multicolumn{3}{|l|}{Anzahl d.\ Sign.\ 6} & \textit{Mittel 1}\,\textit{m}\,\textit{8}\,\textit{s}\,\textit{87} \\
\hline
\end{tabular}
\end{center}

Da des Hofraths Pfaff Garten, und der Gasthof zum Erbprinzen ungefähr in demselben Meridian liegen, so kann man für den Längenunterschied zwischen Seeberg und Helmstedt setzen $1^m\,8^s{,}93$, welches von Paris $34^m\,44^s{,}93$ macht, oder Länge von Ferro $28\degree41'0{,}45''$.

Den 19.~August wurden in des Hofraths Pfaff Garten acht Circum-Meridianhöhen der Sonne beobachtet, welche für die Polhöhe dieses Ortes gaben $52\degree13'37{,}7''$. Acht dergleichen Höhen den 21.~Aug.\ im Gasthofe zum Erbprinzen genommen, gaben die Polhöhe $52\degree13'51{,}5''$. Auf meiner Harzreise im Jahr 1793 fand ich durch einige Winkel-Beobachtungen auf dem Brocken die Breite durch Interpolation für die Hauptkirche in Helmstedt $52\degree12'58''$, für die Länge $28\degree40'10''$, welches für die Methode, n[ämlich]

\newpage
\subparagraph{3.~Wolfenbüttel.}

\begin{center}
\small
\begin{tabular}{|r|c|c|c|}
\hline
\multicolumn{4}{|c|}{\textit{1803}} \\
\hline
& \textit{Mittl.\ Zeit} & \textit{Mittl.\ Zeit} & \textit{Länge in Zeit} \\
& \textit{auf Seeberg} & \textit{in Wolfenbüttel} & \textit{westl.\ v.\ Seeberg} \\
\hline
25.~August & $8^h\,59^m\,41^s{,}87$ & $8^h\,58^m\,54^s{,}58$ & $47^s{,}29$ \\
\quad\, 9 & \quad 41\,,6 & \quad 9\,,53\,,9 & 47\,,7 \\
\quad 19 & \quad 41\,,5 & \quad 18\,,54\,,0 & 47\,,5 \\
\quad 29 & \quad 41\,,8 & \quad 28\,,54\,,1 & 47\,,7 \\
\quad 39 & \quad 41\,,9 & \quad 38\,,54\,,2 & 47\,,7 \\
\quad 49 & \quad 41\,,9 & \quad 48\,,54\,,6 & 47\,,3 \\
\quad 59 & \quad 42\,,1 & \quad 58\,,54\,,3 & 47\,,8 \\
\hline
\multicolumn{3}{|l|}{Anzahl d.\ Sign.\ 7} & \textit{Mittel 47{,}55} \\
\hline
\end{tabular}
\end{center}

welches von Paris östliche Länge $32^m\,47^s{,}55$ und die Länge von Ferro gibt $28\degree11'52{,}5''$. Zehn Meridianhöhen der Sonne in Wolfenbüttel gaben die Polhöhe $52\degree9'29''$. Meine Bestimmung von Wolfenbüttel vom Brocken aus und durch genommene Winkel interpolirt, gibt Breite $52\degree8'44''$, Länge $28\degree11'39''$.\,\,)

\paragraph{(II) [Beobachtungen von Planeten].}

\textsc{Gauss an Olbers.}

(\textsc{Wilhelm Olbers}, Sein Leben und seine Werke, III, 1900, S.~201--202, 234, 236, 281.]

\textit{Braunschweig, 7.~September 1804.} --- [Beobachtungen der Ceres.]

Endlich habe ich auch zum erstenmale das Vergnügen gehabt, unsre Ceres hier in Braunschweig zu beobachten. Den 29.~August ging sie um $32$~Ceti um $11^h\,18^m\,2^s$ nach einem Mittel aus 3 gut harmonirenden Beobb.\ $23{,}5''$ vor und war um $11^h\,31^m\,11^s$ um $1'\,3''$ nördlicher. Daraus finde ich

\begin{center}
\begin{tabular}{l|c|c}
& \textit{Beobachtung} & \textit{Differenz von d.\ Ephem.} \\
\hline
$11^h\,18^m\,2^s$ t.\,m. & AR app.\ $15\degree\,0'\,3''$ & $+3'\,1{,}5''$ \\
& Decl.\ $10\degree\,37'\,56''$ & $+32{,}5$
\end{tabular}
\end{center}

Vorgestern habe ich die Ceres wieder beob[achtet]; allein die Beobb., welche unterbrochen wurden, stimmten schlecht. Nach einer folgte sie

\begin{center}
\begin{tabular}{l@{\quad}l}
um 11\,$^h$1\,$^m$ auf 28 Ceti 22$''$ und ging vor 30 voraus 78$''$ \\
num der 2$^{\text{ten}}$ \quad 11\,$^h$16\,$^m$ auf 28 Ceti 19$''$ und ging voraus vor 30: 81$''$.
\end{tabular}
\end{center}

Decl.\ aus beiden Beobb.\ im Mittel $10\degree36'52''$. Doch ist diese noch zweifelhafter, da der Declinationsunterschied gerade sehr unvortheilhaft für den Durchmesser des Gesichtsfeldes war.\,\,\ldots\ldots

Ich selbst habe die Juno noch am 16., 18., 19. beobachtet und mit dem * 8.\ Größe, der nördl[ich] folgte, verglichen. Die Beobachtungen sind aber alle nicht sehr zuverlässig, besonders die Declination

\begin{center}
\small
\begin{tabular}{c|c|c|c}
& \textit{M[ittlere] Z[eit]]} & \textit{AR} & \textit{Decl.} \\
\hline
Nov.\ 16. & $6^h\,39^m\,56^s$ & $356\degree\,17'\,58''$ & $10\degree\,59'\,25''$ \\
\quad 18. & 6\quad 21\quad\, 6 & 356\quad 33\quad 56 & 10\quad 58\quad 41 \\
\quad 19. & 6\quad\, 4\quad 43 & 356\quad 42\quad 47 & 10\quad 58\quad 34 \quad\text{[wohl zu groß]}
\end{tabular}
\end{center}
[u[nd] nur auf Eine Beobachtung gegründet.\,\,\ldots\ldots]

\textit{Braunschweig, 7.~December 1804.} --- [Beobachtung der Ceres.]

Ceres habe ich am 4.\ Dec[ember] gleichfalls beobachtet:
\[
7^h\,10^m\quad \alpha = 2\degree48'\,9'' = 109\degree\,57'\,45''.
\]
Die Ephemeride gibt also die AR 6' zu klein, die Decl.\ 32' zu groß.\,\,\ldots

\textit{Braunschweig, 7.~December 1804.} --- [Beobachtung der Juno.]

Seit meinem letzten Briefe habe ich die Juno nicht wieder selbst gesehen\ldots

am 4.\ Dec.\ die * einmal wieder selbst beobachtet und mit 17 Bode und dem darauf folgenden * 7.\ 8.\ Größe wiederholt verglichen. Die Position dieser Sterne habe ich selbst aus der Hist.\ Cél.\ abgeleitet und damit gefunden:

\begin{center}
Dec.\ 4.\quad $6^h\,43^m$\quad M.Z.\quad AR\quad $359\degree\,42'\,17''$\quad D.\quad $10\degree\,16'\,3''$,
\end{center}

Die Beobachtungen stimmten nicht schlecht, doch fällt der etwanige Fehler in der Größe des Gesichtsfeldes ganz auf die Declination.\,\,\ldots\ldots

\textit{Braunschweig, 12.~Februar 1805.} --- [Beobachtungen der Juno.]

Ich selbst habe die Juno nur noch zweimal beobachtet, am 26.\ Januar und 1.\ Febr[uar], am letztern Tage eigentlich nur sie gesehen, da 240 Ceti Bode, womit ich sie ein paarmal verglich, zu nördlich stand, um eine brauchbare Position daraus abzuleiten.

Die Beob.\ vom 26.\ Januar ist:
\[
1805\text{ Jan.\ }26.\quad 7^h\,1^m\,47^s\quad 28\degree\,43'\,5''\quad 2\degree\,50'\,53''.
\]
Ich habe diese Beob., die sich auf 5 Vergl[eichungen] mit einem * 8.\ Gr[öße] $17\degree58'$, $3\degree1'$ gründet, noch nicht mit den Elementen verglichen.\,\,\ldots\ldots

\subsubsection*{Bemerkungen}

\textit{Zu I.} --- Der bekannte und einflußreiche Direktor der Seeberger Sternwarte v.\ Zach erhielt vom Könige von Preußen den Auftrag zu einer Gradmessung, die sich über ganz Europa erstrecken sollte, und hatte bereits die Ausführung derselben begonnen, als der Krieg\ldots

\textit{Man sieht leicht, daß} die Methode, welche Gauss für die Längenbestimmung von Braunschweig anwendete, im wesentlichen auf der Vergleichung der mittleren Zeit beruhte. Die Signale dienten dazu, die Uhren zu vergleichen, und die Beobachtungen der Sonnenhöhen gaben die Zeit.

\textit{Die Beobachtungen der Planeten} (II) sind von Gauss selbst angestellt und zeigen die Sorgfalt, mit der er auch die kleinsten Details verfolgte.

\newpage
% ============================================
% SUBSECTION: VERSCHIEDENE BEOBACHTUNGEN IN GOTTINGEN
% ============================================
\subsection*{[Verschiedene Beobachtungen in Göttingen]}

\subsubsection*{Nachlaß, Veröffentlichungen und Briefwechsel}

\textit{Tagebuch der astronomischen Beobachtungen auf der königl.\ Sternwarte zu Göttingen.} (Erster Band; Angefangen November 1808.)

\paragraph{[1.~Seite] 12, 1808 December 18.}
[Austritt des I.\ 4-Trabanten, beobachtet von Herrn Professor Gauss mit 10~f.\ Reflector, schwächste Vergr.

Sternzeit $22^h\,31^m\,38^s{,}0$

$22^h\,29^m\,58^s$ nach Shelton\footnote{Shelton, englischer Uhrmacher.} ist [etwas]; M[ittl.] Z[eit] $= 4^h\,42^m\,44{,}7$\footnote{Vgl.\ Tagebuch No.~130, Werke~XI, 1, S.~565.}.]

\paragraph{[2.~Seite] 55, 1809 September 4.}
Der Eintritt von $\delta$~Geminorum geschah unter dem Horizont; der Austritt wurde beobachtet am dunkeln aber sichtbaren Mondsrande

\begin{center}
\small
\begin{tabular}{l|l|l|l}
& $*$-Zeit & Uhrzeit & M.Z. \\
\hline
von Herr Dr.\ Schumacher mit Short & $0^h\,28^m\,13{,}93$;\quad $0^h\,30^m\,0{,}48$ & $= 13^h\,35^m\,8{,}56$ \\
von mir mit Herschel & \quad 28\quad 13\,,5\quad\quad 30\quad 0\,,68 & = 13\quad 35\quad 8\,,8 \\
von H[errn] H[arding] & \quad 28\quad 13\,,0\quad\quad 30\quad 0\,,41 & = 13\quad 35\quad 9\,,3
\end{tabular}
\end{center}

\paragraph{[3.~Seite] 58, 1809 October 3.}
Bei zieml[ich] heitrer Luft, wiewohl späterhin ein starker Nebel sich verbreitete, wurde die Vesta viermal mit einem Sterne der Histoire Céleste a, p.~253, 1797 Febr.\ 17 verglichen; ein anderer anonymer Stern, b, [der] etwa 15' nördlicher stehen mochte, wurde mitbeobachtet. Vesta hatte die 8.\ Größe.

\begin{center}
\small
\begin{tabular}{c|c|c|c|c|c|c}
& \multicolumn{2}{c|}{S[tern] a} & \multicolumn{2}{c|}{S[tern] b} & \multicolumn{2}{c}{Vesta} \\
\hline
& AR & Decl. & AR & Decl. & AR & Decl. \\
\hline
$2^h\,1^m\,53^s$ & $4^h\,37^m$ & S[üdl.] & $4^h\,41^m\,7^s{,}72$ & $5^h\,58^m$ & $8^h\,53^m\,5^s$ & $2\degree\,57'\,20''$ \\
13 & 2 & 16 & 1 & 15 & 50 & 18\quad 58\,,5 \quad 17\quad 51 \quad 19\quad 38 \\
29 & 19 & 30 & 36 & S. & 32 & 16 \quad 33\quad 17 \quad 32\quad 41 \quad 35\quad 38 \\
40 & 25\,,5 & 43 & 23 & 43 & 14 & 46 \quad 11 \quad 44\quad 52 \quad 47\quad 21 \\
55 & \quad\,\,it & 56 & 53\,,5 & N. & 57 & 52 \quad 59\quad 46 \quad --- \quad ---
\end{tabular}
\end{center}

Reduction des Sterns a aus der Histoire Céleste\ldots\ [$\alpha$] appar.\ $104\degree\,29'\,35{,}82$\quad [$\delta$ app.]\ $19\degree\,50'\,17{,}22$.

Es ging also $\delta$ dem $\alpha$ vor: und war nördlicher als $\alpha$:

\begin{center}
\begin{tabular}{l@{\quad}l|l@{\quad}l}
Uhrzeit & $2^h\,27^m\,25^s{,}85$ & [um] $17^h\,49^m\,39^s$ & [Uhrzeit] $2^h\,9^m\,22^s{,}50$ \\
$*$-Zeit & 2\quad 7\quad 32\,,0 & [$*$-Zeit] & 2\quad 14\quad 28\,,5 \\
M.Z. & 13\quad 14\quad 27
\end{tabular}
\end{center}

\begin{center}
\begin{tabular}{c|c}
$[\alpha =]\quad 104\degree\,1'\,35{,}56$ & $[\delta =]\quad 19\degree\,50'\,28{,}59$
\end{tabular}
\end{center}

Mira Ceti heute Abend 5.\ 6.\ Größe.

\paragraph{[4.~Seite] 98, 1810 August 23.}
[Austritt Nr.\ 130 Tauri Flamst.\ am dunkeln ($\odot$-)Rande

$23^h\,32^m\,37^s{,}25$\quad Gauss 10~f.\ Herschel

\quad\quad\quad\quad\,\,\quad\quad 37\,,8\quad Harpine 5~f.\ Geffken

bei völlig heiterm Himmel. Sehr gute Beob.

Beim Eintritt am hellen ($\odot$-)Rande zitterte der $\zeta$ zu sehr in den Dünsten des Horizonts, so daß der Stern kurz vor dem Eintritte unsichtbar ward.

Wahrer Austritt in $*$-Zeit $= 23^h\,37^m\,45{,}36$ [G]

Mittl.\ Zeit $= 13^h\,31^m\,10{,}46$ [G]

\quad\quad\quad\quad\quad\quad\quad\quad\quad\quad\quad\quad\,\,\,\,\,\, 10\,,8 [H].\footnote{Vgl.\ Tagebuch No.~141.}

\paragraph{[5.~Seite] 109, 1810 September 18.}
[Eintritt $\alpha$ Tauri am hellen ($\odot$-)Rande $22^h\,19^m\,25^s{,}37$ H[arpine] 5~f.\ Reflect\ldots]

\paragraph{[6.~Seite] 127, 1811 März 1.}
Der Austritt des $\alpha$ Tauri hinterm hellen Mondsrande\footnote{Vgl.\ Tagebuch No.~151.}

$8^h\,17^m\,15^s$ Uhrzeit\quad Sternzeit $8^h\,22^m\,19^s{,}89$

M.Z.\ $9^h\,47^m\,16^s{,}4$.

\paragraph{[7.~Seite] 127, 1811 März 7.}

\begin{center}
\small
\begin{tabular}{l|c|c|c|c}
& \multicolumn{2}{c|}{Sternzeit} & \multicolumn{2}{c}{M.Z.} \\
\hline
Eintritt $\sigma$ Leonis am dunkeln Rande & $10^h\,36^m\,41^s$ & $10^h\,41^m\,49^s{,}52$ & $11^h\,42^m\,43^s{,}38$ \\
Austritt & 11\quad 47\quad\, 7 & [vielleicht $1^s$ zu spät] & 11\quad 52\quad 11\,,52 & 12\quad 52\quad 58\,,2
\end{tabular}
\end{center}

\paragraph{[8.~Seite] 134, 1811 Mai 11.}
Heute Abend bei sehr heiterm Himmel wurde die Vesta aufgesucht: ein Stern 7.\ Gr[öße] 1' östlich von dem Platze, der von Herrn Prof[essor] Harding berechnet war. Es stellte sich heraus, daß dieser Fixstern war, während die Vesta\ldots

Hieraus beiläufig\\
\null\quad\, 9)\quad $10^h\,13^m$ M.Z.: AR $246\degree\,49'\,21''$\\
\null\quad 10)\quad 12\quad 29\quad 8.

\paragraph{[9.~Seite] 165, 1812 Febr.\ 8.}
[Von Mira Ceti hatte Herr Professor Gauss diesen Abend mit seinem Cometensucher gar keine Spur finden können\footnote{Vgl.\ Tagebuch No.~158, Werke~XI, 1, S.~568.}.]

\paragraph{[10.~Seite] 176, 1812 Juni 3.}
Heute Abend bestätigte sich, daß der am 1.\ Juni für Juno gehaltene Stern bloß ein Fixstern gewesen sei, was schon am 2.\ Juni bemerkt war. Jedoch glaube ich die Juno heute Abend wirklich erkannt zu haben, da die Configuration der Sterne sich geändert hatte. Hienach nehme ich folgende geschätzte Positionen der Juno an:

\begin{center}
\begin{tabular}{c|c}
Jun.\ 1, & $283\degree17'\quad -4\degree\,54'$ \\
\quad 2, & 283\quad 8\quad -4\quad 52 \\
\quad 3, & 282\quad 59\quad -4\quad 50.
\end{tabular}
\end{center}

Die Ephemeride gibt also die Rectasc[ension] 14' zu klein und die Declination gut oder eine Kleinigkeit zu weit nördlich.

\textit{Was die\ldots}

\newpage
\subsubsection*{Briefwechsel}

\paragraph{[11.] Gauss an Olbers, Göttingen, 28.~Mai 1811.}

(\textsc{Wilhelm Olbers}, Sein Leben und seine Werke, III, 1900, S.~302, 303.]

Am 27.\ März hatte ich auch gewiss die Juno gesehen. Hier meine Beobachtungen:

\begin{center}
\small
\begin{tabular}{c|c|c|c}
\textit{1811} & \textit{Mittl.\ Z.} & \textit{AR} & \textit{Südl.\ Decl.} \\
\hline
April 22 & $9^h\,51^m\,35^s$ & $216\degree\,41'\,50''$ & $0\degree\,58'\,16''$ \\
\quad 24 & 10\quad 32\quad 55 & 216\quad 17\quad 59 & \phantom{0}0\quad 46\quad 50
\end{tabular}
\end{center}

Bogen's Ephemeride gibt also die gerade Aufsteigung jetzt um 7 Minuten zu groß, die Declinationen aber gut. Meine beobachteten Declinationen sind nicht ganz genau, da die Sterne nicht bequem standen. Heute Abend kommt Juno bis auf 1' einem Stern neunter Größe nahe, der nicht in der Histoire céleste steht und den ich Sie zu bestimmen bitte, da ich vielleicht heute Abend die Juno damit vergleiche. Seine scheinbare Stellung, so gut ich sie habe bestimmen können, ist:

\begin{center}
AR.\ $216\degree\,6'\,59''$.\quad Decl.\ austr.\ $0\degree\,38'\,39''$.
\end{center}

Nahe dabei steht ein Stern\ldots

Den Cometen habe ich wegen Mondschein und schlechtem Wetter gar nicht suchen können. Gestern aber habe ich die Vesta beobachtet, aber nur flüchtig so eben reducirt:

\begin{center}
1811 Mai 11.\quad $10^h\,13^m$ M.Z.\quad $246\degree49'$ AR.\quad $12\degree9'$ süd.\ Ab.
\end{center}

Harpine's Ephemeride weicht also $1\degree18'$, die von Bode nur 22' in AR.\ ab. Der Planet ist übrigens leicht zu sehen, denn er hat 6.--7.\ Größe.

\paragraph{[14.] Gauss an Encke, Göttingen, 1.~December 1816.}

Auf Ihren Aufsatz\footnote{J.~F.\ Encke, \textit{Ueber den Cometen von 1812}.} über den Cometen von 1812 freue ich mich im Voraus. Ich habe im Tagebuch [der Sternwarte] nachgesehen; viel wird es nicht sein, was Sie werden gebrauchen können. Ich habe ihn nur eine Nacht beobachtet, am 8.\ Sept.\ [1812] um $2^h\,45^m$ Sternzeit, wo ich ihn mit zwei Sternen der Hist.\ cél.
\begin{center}
A.\ p.\,52\quad $8^h\,3^m\ldots$
\end{center}

\paragraph{[15.] Gauss an Encke, Göttingen, 3.~März 1817.}

Die Reduction der hiesigen Cometenbeobachtungen von 1812 werden Sie bereits vor längerer Zeit von Hrn.\ von Lindenau erhalten\ldots

\paragraph{[16.] Gauss an Gerling, Göttingen, 29.~Mai 1818.}

Künftig soll, hoffe ich, keine Planetenopposition von mir unbeobachtet bleiben. Für die Vesta war es diesmal zu spät. Hoffentlich haben Sie bereits\ldots

\newpage
\paragraph{[17.] Beobachtungen 1818.}

\begin{center}
\small
\begin{tabular}{c|c|c|c}
\textit{[1818]} & & & \\
Aug.\ 31. & $12^h\,25^m\,8^s{,}7$ & $345\degree\,49'\,45{,}1''$ & $+2\degree\,23'\,2{,}4''$ \\
Sept.\ 1. & --- & --- & $+2\quad 10\quad 22\,,9$ \\
\quad 4. & 12\quad 6\quad 26\,,8 & 345\quad\, 5\quad\, 5\,,8 & +1\quad 32\quad 43\,,4 \\
\quad 10. & --- & --- & +0\quad 41\quad\, 9\,,2 \\
\quad 12. & --- & --- & \phantom{+}0\quad 12\quad 33\,,3 \\
\quad 13. & --- & --- & \phantom{+}0\quad 25\quad 43\,,0 \\
\quad 17. & 11\quad\, 5\quad 37\,,5 & 342\quad 39\quad\, 4\,,8 & $-1\quad 19\quad 29\,,5$.
\end{tabular}
\end{center}

\paragraph{[18.] Gauss an Olbers, Göttingen, 9.~Oktober 1825.}

(\textsc{Wilhelm Olbers}, Sein Leben und seine Werke, II~2, Berlin 1909, S.~428.)

Ich habe ihn [den Cometen 1825~IV] zweimal im Meridian zu beobachten versucht, allein die Beobachtungen taugen nicht viel, denn so hübsch er jetzt dem bloßen Auge vorkommt, so elend wird er schon bei der kleinsten 80[fachen] Vergrösserung des Mer[idian]-Kreises und der allergeringsten Beleuchtung, die ihn ganz auslöscht.

Hier die Beobachtungen

\begin{center}
\begin{tabular}{c|c|c}
\textit{1825} & & \\
Okt.\ 4. & $38\degree\,0'\,50''$ & $-14\degree\,6'\,18''$ \\
\quad 6. & 34\quad 34\,[2] & $-18\quad 16\,\,5$
\end{tabular}
\end{center}

\paragraph{[19.] Gauss an Olbers, Göttingen, 19.~Februar 1826.}

(\textsc{Wilhelm Olbers}, Sein Leben und seine Werke, II~2, Berlin 1909, S.~449.)

\ldots bei meinem Umzuge verkramt ist, noch nicht wieder auffinden können. Wollen Sie indessen vorläufig die Sterne reduciren, so will ich demnächst jene Beobachtungen auch berechnen. Leider wimmeln sie nur von Schreibfehlern, die zum Theil durch Conjecturen werden herausgebracht werden müssen.\,\,\ldots
[Folgen Mittheilungen über Harding's Beobachtungen.]

\paragraph{[20.] Gauss an Schumacher, Göttingen, 27.~September 1835.}

Briefwechsel zwischen Gauss und Schumacher, Altona 1860--65, II, S.~421.

Den Cometen [1835~III Halley] habe ich nach dem Mondschein zum ersten Mal wieder beobachtet:

\begin{center}
\begin{tabular}{c|c|c}
\textit{1835} & & \\
Sept.\ 21. & $11^h\,25^m\,26^s$ M.Z. & Ab[weichung] $= 30\degree\,33'\,55{,}6$ \\
& 11\quad 29\quad 25 & G[erade] A[ufsteigung] $= 93\,\,33\,\,20\,,0$
\end{tabular}
\end{center}
ohne noch bei der Reduction auf die Refraction Rücksicht zu nehmen.

\paragraph{[21.] Gauss an Olbers, 11.~November 1835.}

(\textsc{Wilhelm Olbers}, Sein Leben und seine Werke, II~2, Berlin 1909, S.~628.)

Den Cometen [1835~III Halley] habe ich nur wenige Male beobachten können; in der günstigen Zeit habe ich ihn auch nur einmal gesehen, wo sein Anblick etwas sehr Auffallendes darbot.

\newpage
\paragraph{[22.] Astronomische Nachrichten, Band XV, Nr.~345, S.~167.}

\textit{[Sternbedeckung in Göttingen\footnote{Vgl.\ Tagebuch No.~216, Werke~XI, 1, S.~572.}].} --- Herr Hofrath Gauss hat mir folgende von ihm beobachtete Bedeckung des Sterns 576~Arietis mitgetheilt:

\begin{center}
1837 December 9.\quad Eintritt $5^h\,27^m\,37{,}0$ M.Z.\\
\quad\quad\quad\quad\quad\quad\quad\quad Austritt\quad 5\quad 20\quad 13\,,0.
\end{center}

Den Eintritt hält er für sehr gut, den Austritt für unzuverlässig.

\textit{Astronomische Nachrichten, Band XXI, Nr.~494, S.~221.}

\paragraph{[23.] (Gauss an Schumacher) Göttingen 1844 Januar 2.}

Den Cometen [1843~III] habe ich am 1[sten] December zuerst gesehen und an mehreren folgenden Abenden am Kreismikrometer beobachtet; die Beobachtungen sind aber noch nicht reducirt, da ich die Durchmesser der Mikrometerringe noch nicht scharf bestimmt habe. Bloß die letzte vom 13[ten] December kann ich\ldots

\paragraph{[24.] Gauss an Schumacher, Göttingen, 10.~Januar 1844.}

Briefwechsel zwischen Gauss und Schumacher, Altona 1860--65, IV.\ S.~208, V.\ S.~228, 247.

Erst gestern Abend \ldots\ habe ich den Cometen [1843~III] wiedergesehen, nachdem es fast volle vier Wochen ununterbrochen trübe gewesen war. Aber auch gestern Abend war die betreffende Stelle nur zuweilen zwischen Wolken sichtbar. Ich konnte nur eine einzige Vergleichung mit einem Stern 9.\ Größe machen, der dem Cometen etwa 8 Minuten südlich folgte. Für Declination war der Durchgang unbrauchbar: die Rectascension ergab sich

\begin{center}
1844 Jan.\ 9.\quad $8^h\,57^m\,19^s$\quad AR.\ $= 77\degree\,5'\,30{,}8''$
\end{center}
etwa 14 Minuten mehr als die Elemente. Nach einem vollen Monat Zwischenzeit nach den letzten den Elementen zum Gr[öße]

\newpage
\subsubsection*{Bemerkungen}

In dem im Nachlaß befindlichen Tagebuch der Sternwarte, das in zwei Bänden die Zeit vom 3.\ November 1808 bis zum 8.\ Oktober 1818 und vom 30.\ Oktober 1820 bis zum 15.\ April 1822, jedoch mit Unterbrechungen, umfaßt, sind anfangs ziemlich regelmässig, später mehr vereinzelt, die auf der Sternwarte zu Göttingen angestellten Beobachtungen, häufig nebst ihrer Reduction, eingetragen. Die Beobachtungen der ersten Zeit sind vorwiegend von Harding, im besonderen am Mauerquadranten, in der späteren vorwiegend von Gauss angestellt, teilweise auch von Gauss' Schülern und Gästen, Reduziert und veröffentlicht ist davon nur ein Teil. In Nr.\ 1--12 sind aus dem Tagebuch diejenigen noch unveröffentlichten Beobachtungen von Gauss wiedergegeben, die sich auf besondere astronomische Ereignisse oder Objekte beziehen.

\newpage
% ============================================
% SUBSECTION: BEOBACHTUNGEN MIT REICHENBACHSCHEN KREISE
% ============================================
\subsection*{Beobachtungen mit einem 12-zölligen Reichenbachschen Kreise zur Bestimmung der Polhöhe der [alten] Göttinger Sternwarte}

\textit{Monatliche Correspondenz, Band XXVIII, S.~481--484, 1813 Mai.}

\textit{Aus einem Schreiben des Hrn.\ Prof.\ Gauss. Göttingen, 18.~Mai 1813.}

Sie wissen, theuerster Freund, welche zufällige Hindernisse den Anfang der astronomischen Beobachtungen mit dem herrlichen Reichenbachschen Kreise, welchen unsere Sternwarte gegen Ende des vorigen Jahres erhielt, bis zur Mitte des März verzögert haben. Ich habe jetzt das Vergnügen, Ihnen die Erstlinge meiner Beobachtungen mitzutheilen, welche die Bestimmung der Polhöhe der Sternwarte zum Zweck haben. Auf der Südseite des Meridians, wo ich dem Kreise eine eben so feste Aufstellung wie er jetzt für die Nordseite hat, noch nicht habe \ldots

\begin{center}
\small
\begin{tabular}{|l|c|c|c|c|}
\hline
& \textit{Zenith-Distanz} & \textit{Refraction} & \textit{Scheinb.} & \textit{Polhöhe} \\
& & & \textit{Zenith-Dist.} & \\
\hline
März 20. & $40\degree\,9'\,25{,}48''$ & $1\degree\,41'\,20{,}34''$ $(19{,}63)$ & $(51\degree\,31'\,54{,}86'')$ & $(54{,}515)$ \\
\quad 22. & 25\,,19 & 20\,,92 $(20{,}24)$ & 55\,,73 $(55{,}05)$ \\
\quad 26. & 24\,,68 & 22\,,09 $(21{,}48)$ & 57\,,41 $(56{,}80)$ \\
\quad 31. & 27\,,33 & 23\,,58 $(23{,}02)$ & 56\,,25 $(55{,}69)$ \\
April\,\,3. & 27\,,91 & 24\,,46 $(23{,}93)$ & 56\,,55 $(56{,}02)$ \\
\quad\,\, 7. & 30\,,44 & 25\,,64 $(25{,}16)$ & 55\,,20 $(54{,}72)$ \\
\quad\,\, 8. & 28\,,59 & 25\,,92 $(25{,}46)$ & 57\,,33 $(56{,}87)$ \\
\hline
\multicolumn{4}{|l|}{\textit{Mittel aus 136 Beobachtungen}} & \textit{51\degree31'}\,\textit{56{,}20''} $(55{,}64)$ \\
\hline
\end{tabular}
\end{center}

Hieraus also:

\begin{center}
\small
\begin{tabular}{|l|c|c|}
\hline
& \textit{Polhöhe} & \\
\hline
März 20. & $51\degree\,31'\,54{,}86''$ & $(54{,}15)$ \\
\quad 22. & 55\,,73 & $(55{,}05)$ \\
\quad 26. & 57\,,41 & $(56{,}80)$ \\
\quad 31. & 56\,,25 & $(55{,}69)$ \\
April 3. & 56\,,55 & $(56{,}02)$ \\
\quad 7. & 55\,,20 & $(54{,}72)$ \\
\quad 8. & 57\,,33 & $(56{,}87)$ \\
\hline
\textit{Mittel} & \textit{51\degree31'}\,\textit{56{,}20''} & $(55{,}64)$ \\
\hline
\end{tabular}
\end{center}

Bei Berechnung dieser Beobachtungen liegt die Polar-Distanz $\alpha$ im kleinen Bär für den Anfang von 1812
\[
= 1\degree\,41'\,41{,}574''
\]
zum Grunde\footnote{Pond's Bestimmung.}). Herr Pond hat sie mit dem neuen Troughtonschen Kreise in Greenwich für dieselbe Epoche aus Sommerbeobachtungen $= 1\degree\,41'\,41{,}60''$ aus Winterbeobachtungen $= 1\quad 41\quad 41\,,00$ gefunden; die\ldots

Sternwarte $= 51\degree\,31'\,56{,}04''$ oder $= 51\degree\,31'\,55{,}560$ $(51\degree\,31'\,55{,}48'')$ je nachdem wir für die Declination des Polarsterns blos Herrn von Zach's Bestimmung, oder das Mittel von vier Astronomen annehmen. Mayer's Bestimmung dieser Polhöhe wäre demnach nur $2''$ oder $1{,}76''$ zu klein.

Ich habe in diese Reihe die Beobachtungen von drei andern Abenden nicht mit aufgenommen, weil dies meine allerersten Versuche waren; inzwischen stimmen auch diese auf das schönste überein, es sind folgende:

\begin{center}
\small
\begin{tabular}{|l|l|}
\hline
1813 März 14. & Vier Zenith-Distanzen des Polarsterns ausser\\
& der Culmination geben Polhöhe \ldots\ldots\quad $51\degree\,31'\,56{,}57''$ \\
\quad\quad\,\, 15. & Zenith-Distanz von $\delta$~Cephei in der untern\\
& Culmination 4 Beobachtungen \ldots\ldots\quad 56\,,14 \\
\quad\quad\,\, 17. & Zenith-Distanz von $\delta$~Cephei in der obern\\
& Culmination 4 Beobachtungen \ldots\ldots\quad 55\,,84 \\
\hline
\end{tabular}
\end{center}

\newpage
% ============================================
% SUBSECTION: UEBER DAS KREISMIKROMETER
% ============================================
\subsection*{Über das Kreismikrometer}

\subsubsection*{Briefwechsel}

(\textsc{Wilhelm Olbers}, Sein Leben und seine Werke, III, Berlin 1900, S.~402, 409 und II~2, Berlin 1909, S.~69.]

\paragraph{Gauss an Olbers, Göttingen, 5.~Januar 1808.}

Mein Verfahren zur Reduction der Kreismikr[ometer-]Beobacht[ungen] mit Rücksicht auf Refraction gründet sich darauf, daß in dem kreisförmigen Gesichtsfelde sich vom Himmel ein Ausschnitt zeigt, der als elliptisch angesehen werden kann: die große Axe dieser Ellipse liegt in dem Vertikalkreise; diese Voraussetzung weicht nur ganz unmerklich von der Wahrheit ab. Es sei

\begin{center}
\begin{tabular}{rl}
$2r$ & wahrer Halbmesser des Gesichtsfeldes, \\
$2kr$ & große (vertikale) Axe der Ellipse, \\
$2hr$ & kleine (horizontale) Axe der Ellipse,
\end{tabular}
\end{center}

so werden $h$, $k$ von der Höhe abhängige Faktoren, etwas weniges größer als 1 sein. Setzt man für

\begin{align*}
\gamma' - \delta &= D + kr\sin b\sqrt{1 - A^2\sin p^2} = D + kr\sin b\sqrt{1 - AA\sin p^2}\\
\gamma - \delta &= D - kr\sin b\sqrt{1 - A^2\sin p^2}
\end{align*}
$A$ kann man für alle Höhen als constant betrachten; ich finde
\[
A^2\cos p^2 \sqrt{1 - A\sin p^2} = \cos\psi.
\]
Es ist also nur eine Tafel für $\psi$ und $A$ erforderlich; hier eine solche, flüchtig berechnet:

\begin{center}
\small
\begin{tabular}{r|cc||r|cc||r|cc||r|cc}
$z$ & $\log k$ & $\log A$ & $z$ & $\log k$ & $\log A$ & $z$ & $\log k$ & $\log A$ & $z$ & $\log k$ & $\log A$ \\
\hline
$0\degree$ & 0\,,00000 & $-\infty$ & $22\degree$ & 0\,,00002 & 7\,,982 & $44\degree$ & 0\,,00011 & 8\,,360 & $66\degree$ & 0\,,00061 & 8\,,723 \\
1 & 6\,,618 & & 23 & 2 & 8\,,003 & 45 & 12 & 375 & 67 & 73\,[3] & 743 \\
2 & 6\,,919 & & 24 & 2 & 024 & 46 & 13 & 390 & 68 & 81 & 764 \\
3 & 7\,,095 & & 25 & 3 & 044 & 47 & 14 & 405 & 69 & & 786 \\
4 & 7\,,22 & & 26 & 3 & 064 & 48 & 16 & 419 & 70 & & 809
\end{tabular}
\end{center}

\begin{align}
\log A &= 0{,}00012 \text{ von } z = 0\degree \text{ bis } 82\degree \notag\\
\log k &= 0{,}00011 \text{ von } z = 83\degree \text{ bis } 85\degree \notag\\
\log k &= 0{,}00010 \text{ für } z = 86\degree \notag\\
\log k &= 0{,}00009 \text{ für } z = 87\degree. \notag
\end{align}

\paragraph{Gauss an Olbers, Göttingen, 11.~Februar 1808.}

Bessel's Methode\footnote{Vgl.\ Werke~VIII, S.~373.}, Kreismikrometerbeobachtungen mit Rücksicht auf Refraction zu reduciren, kenne ich nicht, und kann also in Rücksicht auf die Schärfe zwischen derselben und der meinigen [keine Vergleichung] anstellen. Bei der meinigen liegen zwei nur näherungsweise richtige Voraussetzungen zum Grunde, [1.] daß das kreisförmige Gesichtsfeld durch die Refraction elliptisch wird (wenn ich mir diese Phrase erlauben darf), 2.\ daß das Verhältnis der Dimensionen der Ellipse zu denen des Kreises unabhängig von der Größe des Gesichtsfeldes ist. Der letzten kann man überhoben sein, wenn man sich für sein Fernrohr eine besondere Tabelle construirt, wie die meinige, die eigentlich für ein unendlich kleines Gesichtsfeld berechnet ist.

Auch das Gesichtsfeld ist als genau bekannt betrachtet. Aber auch ohne diese Umstände zu berücksichtigen, wird es doch viel genauer sein, die Beobachtungen nach jenen Gewichten zu berechnen, als wenn man alle Resultate als gleich genau ansieht. Ich habe dies bei den Cometenbeobb.\ gethan. Die Rechnung ist höchst einfach, wenn man sich meiner kleinen Logarithmentafeln bedient; ich gehe für die erste Formel mit der Differenz der Logarithmen von $\cos\varphi$ u[nd] $\cos\varphi'$ in die Kolumne $A$ ein und subtrahire $B$ allemal von dem kleinern Logarithmen von $2\cos\varphi'$ oder $2\cos\varphi''$. Ebenso bei der andern Formel. Am 10.\ Februar waren z.B.\ meine einzelnen Resultate\footnote{A.\ N.\ Band I, Nr.\ 14, S.\ 216.}

\begin{center}
\small
\begin{tabular}{c|c|c|c|c}
& \multicolumn{2}{c|}{für Decl.-Unterschied} & \multicolumn{2}{c}{für AR} \\
\hline
& & Gewicht & & Gewicht \\
\hline
13'\,55{,}5 & & 0\,,666 & 47'\,52{,}5 & 0\,,294
\end{tabular}
\end{center}

\newpage
\subsubsection*{Bemerkungen}

Dann sind die Gleichungen der Ellipse und der Sehne
\[
\alpha x^2(1 - A^2) + y^2 = a^2,\quad y\cos\alpha = x\sin\alpha + d,
\]
also die Schnittpunkte beider entsprechend dem Ein- und Austritt des Sterns
\[
x^2(1 - A^2\cos\alpha^2) = -d\sin\alpha \pm \sqrt{\{(1 - A^2)(a^2(1 - A^2\cos\alpha^2) - d^2)\}}
\]
und hieraus folgt für die Länge der Sehne
\[
\frac{v^2}{4} = \frac{a^2}{\cos^2\psi} - \frac{(1 - A^2)\{a^2(1 - A^2\cos\alpha^2) - d^2\}}{\cos^2\psi(1 - A^2\cos\alpha^2)^2}.
\]

Führt man den Winkel $\vartheta$ ein, indem man setzt
\[
d^2 = a^2\sin^2\vartheta(1 - A^2\cos\alpha^2),
\]
so wird
\[
\frac{v^2}{4} = \frac{a^2(1 - A^2)\cos\vartheta^2}{1 - A^2\cos\alpha^2}.
\]

Da
\[
d = \delta - D,\quad p\alpha = (\beta' - \beta)\cos\delta,\quad p = 90\degree - \alpha,
\]
so folgen die beiden Gaußschen Gleichungen I und II aus den beiden vorstehenden. Ferner hat man für die Abszisse der Sehnenmitte
\[
\frac{1}{2}(y + y') = -\frac{d\sin\alpha}{1 - A^2\cos\alpha^2}.
\]

Die Abszisse des Fußpunktes des vom Mittelpunkt auf die Sehne gef[ällten] Lotes ist\ldots

\newpage
% ============================================
% SUBSECTION: UEBER DAS HELIOMETER
% ============================================
\subsection*{[Über das Heliometer]}

\subsubsection*{Briefwechsel}

\paragraph{Gauss an Gerling, Göttingen, 17.~Junius 1814.}

Die Art, wie unser Heliometer das Repetiren bewirkt, ist äusserst einfach. Die Hauptsache ist, daß beide Objectivhälften für sich beweglich sind; nur die eine Hälfte gibt aber feine Ablesung, diese heisse $A$, die andere $B$. Jedes Halbobjectiv macht von jedem Objecte ein Bild. Die Objecte seien 1, 2, so habe ich vier Bilder $A_1$, $A_2$, $B_1$, $B_2$. Ich bringe erst (gleichviel auf welche Art) $A_1$ mit $B_2$ in Berührung und lese ab; dann bringe ich durch Bewegung von $A$, $A_2$ mit $B_1$ in Berührung und lese wieder ab, so gibt der Unterschied den doppelten Winkel; jetzt lasse ich $A$ ruhen und bringe durch Bewegung von $B$ wieder $A_1$ mit $B_2$ in Berührung, dann wieder durch Bewegung von $A$, $A_2$ mit $B_1$ und lese wieder ab; auf diese Weise erhalte ich das vierfache etc.; beide Objectivhälften bewegen sich hiebei immer wechselsweise in einerlei Sinn und die Repetition kann so weit getrieben werden als es die Größe des [Feldes zuläßt.]

Man kann noch durch eine kleine Mechanik am Heliometer den Positionswinkel unmittelbar messen, ein Vortheil, der bei der gewöhnlichen Aufstellung wegfallt. Die vier Vergrösserungen schätze ich auf 50, 70, 100, 150mal. Seit einigen Tagen beobachte ich einen sehr kleinen Sonnenfleck (nur ein Paar Secunden groß), ich hoffe, dass Nicolai ihn zugleich am P[assagen-]I[nstrument] beobachten wird; dann werden sich gute Resultate ergeben. Ich dachte ihn zugleich am Quadranten zu beobachten, allein es ist nur an Einem Mittage heiter gewesen, und da zeigte das elende Fernrohr desselben den Fleck beinahe gar nicht; auch ist die Zeit beinahe zu kurz für Einen Beobachter, beide Ränder und den Flecken zu nehmen und sechsfach abzulesen.\,\,\ldots\ldots

(\textsc{Wilhelm Olbers}, Sein Leben und seine Werke, III, Berlin 1900, S.~555, 586, 592.]

\paragraph{Gauss an Olbers, Göttingen, 7.~Juli 1814.}

Mein Heliometer gibt die Rotation der Objektive nur in Graden an; ein geübtes Auge fehlt aber beim Schätzen wohl nicht leicht über $3'$, und da\ldots

\newpage
\subsubsection*{Bemerkungen}

\ldots die Scale abzuleiten, so kommt der Umstand in Betracht, daß nach aller Strenge die Verschiebung des Objectivs dem Winkel nicht proportional ist. Soll man nun die Tangente der Hälfte, den Sinus der Hälfte oder noch etwas anderes nehmen? Das soll billig theoretisch entschieden werden, aber hier eben ist nun die Schwierigkeit. Wohin soll man eigentlich das Bild setzen von Strahlen, die nicht parallel mit der Axe auf ein Objectiv fallen? Ich denke mir einen Hauptstrahl, der auf die Mitte des Objectivs $C$ (geneigt gegen die Axe) fällt. Die Ebene durch den Strahl und die Axe gehe durch $CD$, und $D$ sei dem $C$ unendlich nahe, so kann man allerdings den Punkt bestimmen, wo die auf $C$ und $D$ fallenden Parallelstrahlen sich schneiden, und wenn das Objectiv nach richtigen Principien gearbeitet ist, d.\,i.\ so, daß die sogenannte Abweichung wegen der Kugelgestalt gehoben ist, so werden auch alle übrigen Parallelstrahlen, die in derselben [Ebene liegen, sich in demselben Punkte schneiden.]

\ldots zusinnen. Bisher habe ich sie aber erst im Groben gemacht, da in dem neuen Ocular erst Spinnfaden eingezogen werden müssen, woran es bisher gefehlt hat.\,\,\ldots\ Aber auch bei der frühern groben Berichtigung fand ich die Venus sogleich am Tage. Sehen Sie zur Probe die Messung des Venusdurchmessers am 16.\ Mai Nachmittags

\begin{center}
\small
\begin{tabular}{c|c|c|c}
Abfahrungspunkt & Partes & & \\
\hline
50 fach & 56\,,4 & & \\
2 fach & 50\,, & 13\,,6 & $21{,}40 = 12{,}08$ \\
4 fach & 49\,, & \,71\,,1 & $21{,}32 = 12{,}04$ \\
6 fach & 49\,, & \,30\,,0 & $21{,}07 = 11{,}90$ \\
8 fach & 48\,, & \,88\,,6 & $20{,}97 = 11{,}84$ \\
10 fach & 48\,, & \,47\,,6 & $20{,}88 = 11{,}79$
\end{tabular}
\end{center}

Die theoretische Schwierigkeit, wovon ich Ihnen in meinem letzten Briefe schrieb, hat insofern eine praktische Wichtigkeit, als ich gern den genauen Werth der Scale aus einem so großen Winkel, als das Instrument nur faßt, bestimmen möchte. Ich denke dabei in Zukunft etwa folgendes Verfahren anzuwenden. Ich lasse etwa 6 Stangen $A$, $B$, $C$ etc.\ ausstecken

\begin{center}
$A\quad B\quad C\quad D\quad E\quad F$
\end{center}

so daß $AB$, $BC$, $CD$ etc.\ ungefähr unter sich gleich sind und etwa einen Winkel von $1\frac{1}{2}$ [Grad bilden.]

\newpage
% ============================================
% SUBSECTION: BERICHTE UND ANTRAEGE
% ============================================
\subsection*{[Berichte und Anträge zur Anschaffung von Instrumenten für die Göttinger Sternwarte]}

\textit{[Aus den Akten des Kuratoriums der Universität Göttingen.]}

\subsubsection*{[1.] Unterthänigstes Pro\,\,M[emoria] die Anschaffung der Instrumente für die neue Sternwarte in Göttingen betreffend.}

Da der Bau des Hauptgebäudes der neuen Sternwarte im Wesentlichen jetzt so gut wie vollendet ist, und nach dem Beschluß des hohen königl.\ Curatorium der Bau der Seitenflügel im nächsten Frühjahr mit Nachdruck betrieben werden wird, so rückt nun die gänzliche Vollendung dieses für die Universität und für die Wissenschaft selbst gleich wichtigen Instituts immer näher. Schon zwölf Jahre sind seit seiner ersten Gründung verflossen. Es ist gewiß, daß die Erwartung des ganzen für die Astronomie sich interessirenden Europa auf die künftige Thätigkeit desselben gerichtet ist. Es ist höchst wünschenswerth, daß diese Thätigkeit sofort nach Vollendung des Baues auf eine würdige Art anfangen könne: allein da\ldots

\ldots auf die von Lilienthal zu erwartenden, als endlich auf die neu anzuschaffenden [Instrumente] ganz bestimmte Rücksicht genommen ist. Da die Hauptinstrumente besonders unter den letztern eine äusserst feste Aufstellung erfordern, so ist für jedes derselben gleich zu Anfang ein besonderes 12~Fuß tiefes Fundament aus festen Quadern gelegt. Freilich hat in dem beträchtlichen seitdem verflossenen Zeitraume die Instrumentalastronomie in mehrern Stücken wichtige Verbesserungen und Abänderungen erhalten, die auch bei jenem Plane einige Modificationen nöthig machen werden: inzwischen werden doch jene Fundamente im Wesentlichen ihre ursprüngliche Bestimmung behalten.

Um nun beurtheilen zu können, was von den neu anzuschaffenden Instrumenten zuerst und am dringendsten nöthig sein wird, muß wenigstens eine allgemeine Uebersicht der Bedürfnisse einer Sternwarte vom ersten Range, wie die hiesige werden soll, vorausgehen, wobe\ldots

\ldots ersteren Güte müssen demnach noch hinzukommen. Zu den allerdringendsten Bedürfnissen gehören indessen dieselben nicht. In der Rotunde ist nicht eher eine Uhr nöthig, als bis das dahin zu setzende Aequatorealinstrument zugleich da ist, so wie in den beiden Abtheilungen der untern Sternwarte die Uhren zugleich mit den daselbst aufzustellenden Instrumenten da sein müssen. Ueberdies hat die Sternwarte in Zukunft aus Lilienthal auch eine Pendeluhr zu erwarten, die wie ich glaube als eine von den Hauptuhren wird dienen können.

\textit{Sehwerkzeuge} d.\,i.\ achromatische Fernröhre und Spiegelteleskope. Die alte Sternwarte besitzt von Instrumenten dieser Classe mehrere recht gute; noch weit mehr wird sie aber in Zukunft aus Lilienthal erhalten, wo sich bekanntlich ein großer Reichthum von Spiegelteleskopen, zum Theil von großen Dimensionen, befindet. In dieser Beziehung bleibt demnach wenig zu wünschen übrig.

\textit{Bewegliche Meßinstrumente.} Aus dies[en]

\newpage
\subsubsection*{[Fortsetzung]}

Dem von Anfang an entworfenen Plane zufolge soll die neue Sternwarte folgende fünf feste Instrumente enthalten, für welche schon im Voraus alle Einrichtungen berechnet und getroffen sind.

\begin{center}
\begin{tabular}{cp{10cm}}
A. & Ein Aequatoreal, in der Rotunde aufzustellen, wofür das Drehdach eingerichtet ist. \\
B. & Ein großes Passageinstrument, unten, in der westlichen Abtheilung. \\
C. & Ein großer Meridiankreis, ebendaselbst. \\
D. & Ein kleineres Passageinstrument, in der östlichen Abtheilung. \\
E. & Der schon auf der alten Sternwarte seit 60 Jahren befindliche Mauerquadrant ebendaselbst.
\end{tabular}
\end{center}

Für die Instrumente B, C, D, E sind die oben erwähnten Fundamente gelegt und die beiden Durchschnitte des Gebäudes sowie die Dachklappen eingerichtet.

Es wird nöthig sein, von diesen Instrumenten wenigstens mit einigen Zügen den Gebrauch zu bezeichnen.

Als Hauptzweck aller großen Sternwarten in wissenschaftlicher Rücksicht muß angesehen werden \textit{die Bewegungen und [Oerter der Fixsterne} zu bestimmen, und zwar die eigentlichen Oerterveränderungen derselben und die jährliche Präcession, Nutation, Aberration, Refraction, Parallaxe.]

\vspace{0.5em}
Aus dieser kurzen Darstellung geht schon hervor, daß die Instrumente B, C als die wichtigsten der ganzen Sternwarte angesehen werden müssen. Diese, wenn sie von der ersten Güte sind, erheben eigentlich die Sternwarte zum ersten Range, und die Beobachtungsregister, welche von den ersten Sternwarten z.\,B.\ der Greenwicher jährlich gedruckt werden, enthalten fast gar nicht anderes als die täglichen, mit den Meridianinstrumenten angestellten Beobachtungen, und auf diesen seit 50 oder 60 Jahren gedruckten Greenwicher Beobachtungen beruhet der größte Theil der heutigen Astronomie. Nicht eher wird sich die Göttinger Sternwarte neben die Greenwicher, Pariser, Mailänder, Palermer und die eben jetzt zur Vollendung kommenden Ofner und Neapler Sternwarten stellen können, als bis sie mit tüchtigen Meridianinstrumenten versehen ist.

Uebrigens muß ich noch bemerken, daß dem Plane nach die Instrumente D, E hauptsächlich für den Unter[richt bestimmt sind.]

\end{document}
